\documentclass[12pt,a4paper]{report}

% Packages
\usepackage[utf8]{inputenc}
\usepackage[T1]{fontenc}
\usepackage{amsmath,amssymb}
\usepackage{graphicx}
\usepackage{booktabs}
\usepackage{hyperref}
\usepackage[margin=1in]{geometry}
\usepackage{caption}
\usepackage{subcaption}
\usepackage{algorithm}
\usepackage{algorithmic}
\usepackage{cite}
\usepackage{xcolor}
\usepackage{listings}
\usepackage{enumitem}

% Custom commands
\newcommand{\ddG}{$\Delta\Delta G$}
\newcommand{\dG}{$\Delta G$}
\newcommand{\PCC}{\text{PCC}}
\newcommand{\todo}[1]{\textcolor{red}{\textbf{[TODO: #1]}}}

% Title
\title{
    \textbf{DeepPEF: Improving Protein Thermodynamic Stability Prediction\\via Structure-Aware Embeddings and Graph Neural Networks}\\[1cm]
    \large M.Sc. Thesis\\
    \large Department of Computer Science
}
\author{[Student Name]}
\date{\today}

\begin{document}

\maketitle

\begin{abstract}
Predicting the thermodynamic effect of amino acid mutations on protein stability ($\Delta\Delta G$) is critical for protein engineering, drug design, and understanding genetic disease. This thesis continues and extends the DeepEF (Deep Protein Energy Function) framework, a graph neural network architecture that combines protein language model embeddings with structural features for $\Delta\Delta G$ prediction. Starting from a baseline Pearson correlation coefficient (PCC) of 0.467 on the MegaScale benchmark, we systematically improve the model through loss function engineering, architectural ablation, and embedding upgrades. Our best single-variable improvement --- $k$-nearest neighbor graph attention --- achieves PCC = 0.510. We present a staged plan to reach PCC $\geq$ 0.70 through structure-aware SaProt embeddings, training optimization, per-protein fine-tuning, and model ensembling. All experiments are conducted on the MegaScale/PNAS dataset (340 training proteins, 28 test proteins, $\sim$50,000 single-point mutations) without access to pretrained checkpoints.
\end{abstract}

\tableofcontents

\chapter{Introduction}
\label{ch:introduction}

\section{Motivation}

Proteins are the molecular machines of life. Their function --- catalysis, signaling, structural support --- depends critically on their three-dimensional structure, which in turn depends on thermodynamic stability. A protein's folding free energy ($\Delta G$) determines whether it adopts its native fold or remains unstructured. When a single amino acid is mutated, the resulting change in stability ($\Delta\Delta G = \Delta G_{\text{mutant}} - \Delta G_{\text{wild-type}}$) can determine whether the protein functions normally, loses activity, or gains pathogenic properties.

Accurate computational prediction of $\Delta\Delta G$ upon mutation has profound applications:
\begin{itemize}
    \item \textbf{Drug design:} Engineering therapeutic proteins with enhanced stability
    \item \textbf{Enzyme engineering:} Designing industrial enzymes that tolerate high temperatures
    \item \textbf{Disease understanding:} Identifying destabilizing mutations that cause genetic disorders
    \item \textbf{Directed evolution:} Guiding laboratory evolution by predicting beneficial mutations
\end{itemize}

Despite decades of research, accurate $\Delta\Delta G$ prediction remains challenging. Physics-based methods (FoldX, Rosetta) capture known biophysical principles but struggle with accuracy beyond 1--2 kcal/mol RMSE. Machine learning approaches have shown promise, but until recently lacked sufficient training data to generalize across protein families.

\section{The MegaScale Breakthrough}

In 2023, Tsuboyama et al.\ published the MegaScale dataset \cite{tsuboyama2023mega}, measuring $\Delta\Delta G$ for $\sim$776,000 mutations across $\sim$580 small protein domains using a high-throughput proteolysis assay. This dataset is 100$\times$ larger than previous resources (ProTherm: $\sim$7,000 entries) and provides unprecedented coverage for training deep learning models.

The availability of MegaScale, combined with advances in protein language models (ESM-2, ProtT5, SaProt) and graph neural networks, creates an opportunity to build highly accurate $\Delta\Delta G$ predictors.

\section{DeepEF: Background}

DeepEF (Deep Protein Energy Function) was originally developed by Shahar Cohen as a graph neural network framework for protein stability prediction. The model:
\begin{enumerate}
    \item Constructs a protein graph from 3D backbone coordinates (N, CA, C, CB atoms)
    \item Encodes residue features using protein language model embeddings (originally ProtT5, 1024-dim)
    \item Processes the graph through parallel GCN (bonded features) and GAT (non-bonded features) branches
    \item Applies Light Attention pooling for sequence-level aggregation
    \item Predicts energy values for folded/unfolded states; $\Delta\Delta G$ = $E_{\text{mutant}} - E_{\text{wild-type}}$
\end{enumerate}

The original DeepEF achieved PCC $\approx$ 0.54 with a pretrained checkpoint on native/decoy discrimination (CASP12 dataset) followed by fine-tuning on MegaScale mutations.

\section{Thesis Objective}

This thesis continues DeepEF with the following constraints and goals:
\begin{itemize}
    \item \textbf{Starting point:} No access to the original pretrained checkpoint --- all training from scratch
    \item \textbf{Dataset:} MegaScale/PNAS filtered subset (340 train proteins, 28 test proteins)
    \item \textbf{Hardware:} Single NVIDIA RTX 4070 (8GB VRAM), WSL2 environment
    \item \textbf{Target:} Improve PCC from 0.467 (baseline) to $\geq$ 0.70 through systematic improvements
    \item \textbf{Approach:} Embedding upgrades (ProtT5 $\rightarrow$ SaProt), architectural ablation, training optimization, and ensembling
\end{itemize}

\section{Contributions}

The main contributions of this thesis are:
\begin{enumerate}
    \item \textbf{Systematic ablation study} identifying which architectural modifications help ($k$-NN GAT: +0.043 PCC) and which hurt (Multi-RBF: --0.081 PCC, edge features: neutral)
    \item \textbf{Loss function engineering} showing Huber + Ranking loss outperforms L1 loss for $\Delta\Delta G$ regression
    \item \textbf{Training pipeline optimization} with cosine LR scheduling, gradient accumulation, and weight decay
    \item \textbf{Embedding upgrade strategy} from ProtT5 (1024-dim, sequence-only) to SaProt (1280-dim, structure-aware), with validated integration pipeline
    \item \textbf{End-to-end reproducible framework} for protein stability prediction without pretrained checkpoints
\end{enumerate}

\section{Thesis Structure}

Chapter~\ref{ch:related_work} reviews related work in protein stability prediction, protein language models, and graph neural networks for proteins. Chapter~\ref{ch:methods} describes the DeepEF architecture, training pipeline, and data processing in detail. Chapter~\ref{ch:experiments} presents our systematic experiments: loss ablation, architecture ablation, and the Phase~2 failure analysis. Chapter~\ref{ch:results} summarizes current results. Chapter~\ref{ch:future_work} presents the staged improvement plan with expected gains and timeline.

\chapter{Related Work}
\label{ch:related_work}

\section{Protein Stability Prediction}

Protein thermodynamic stability prediction has evolved through several paradigms:

\subsection{Physics-Based Methods}

\textbf{FoldX} \cite{schymkowitz2005foldx} uses an empirical force field to estimate $\Delta\Delta G$ from structural features (van der Waals, solvation, hydrogen bonds, entropy). It achieves $\sim$1.0 kcal/mol RMSE on curated datasets but requires high-quality input structures and struggles with backbone flexibility.

\textbf{Rosetta ddg\_monomer} \cite{kellogg2011role} performs Cartesian-space energy minimization after mutation modeling. More accurate than FoldX for buried mutations but computationally expensive ($\sim$10 minutes per mutation).

\subsection{Machine Learning Methods}

\textbf{ThermoMPNN} \cite{dieckhaus2024thermompnn} adapts the ProteinMPNN message-passing architecture for $\Delta\Delta G$ prediction. Pre-trained on inverse folding then fine-tuned on stability data, it achieves PCC $\approx$ 0.75 on MegaScale benchmarks with structure-aware inputs.

\textbf{RaSP} (Rapid Stability Prediction) \cite{blaabjerg2023rasp} trains a 3D CNN on Rosetta-generated $\Delta\Delta G$ landscapes, achieving high throughput with competitive accuracy.

\textbf{ACDC-NN} \cite{benevenuta2021acdc} uses sequence and evolutionary features with a feed-forward neural network.

\textbf{DDGun3D} \cite{montanucci2019ddgun} combines sequence conservation with structural features in a gradient boosting framework.

\section{Protein Language Models}

Pre-trained protein language models encode evolutionary and structural information from millions of protein sequences:

\subsection{ProtT5}

ProtT5 \cite{elnaggar2022prottrans} is a T5-based encoder-decoder model trained on UniRef50 ($\sim$45M sequences). It produces 1024-dimensional per-residue embeddings. Key properties:
\begin{itemize}
    \item Captures evolutionary conservation and co-evolution patterns
    \item Sequence-only: no structural information encoded
    \item Thermostability benchmark: Spearman $\rho \approx 0.65$
\end{itemize}

\subsection{ESM-2}

ESM-2 \cite{lin2023evolutionary} is a transformer-based masked language model (15M to 15B parameters). The 650M variant produces 1280-dimensional per-residue embeddings. Compared to ProtT5:
\begin{itemize}
    \item Trained on UniRef50 with updated architecture (RoPE, no bias, pre-normalization)
    \item Better mutation effect prediction: Spearman $\rho = 0.680$ on thermostability
    \item Sequence-only but with stronger evolutionary signal due to scale
\end{itemize}

\subsection{SaProt}

SaProt (Structure-aware Protein language model) \cite{su2024saprot} extends ESM-2 by incorporating Foldseek 3Di structural tokens:
\begin{itemize}
    \item Input: interleaved amino acid + 3Di tokens (e.g., ``MdEvVpQp'')
    \item 3Di tokens encode local backbone geometry (torsion angles, inter-residue distances)
    \item Same architecture as ESM-2 (650M params, 1280-dim output)
    \item \textbf{Thermostability: Spearman $\rho = 0.724$} --- best among PLMs (+6.5\% over ESM-2)
    \item Ranked \#1 on ProteinGym leaderboard (April 2024)
    \item Requires AlphaFold-predicted structures for 3Di token extraction
\end{itemize}

\begin{table}[h]
\centering
\caption{Protein language model comparison on thermostability prediction}
\label{tab:plm_comparison}
\begin{tabular}{lccc}
\toprule
\textbf{Model} & \textbf{Params} & \textbf{Embedding Dim} & \textbf{Spearman $\rho$} \\
\midrule
ProtT5-XL & 3B & 1024 & $\sim$0.65 \\
ESM-2 & 650M & 1280 & 0.680 \\
\textbf{SaProt} & \textbf{650M} & \textbf{1280} & \textbf{0.724} \\
\bottomrule
\end{tabular}
\end{table}

\section{Graph Neural Networks for Proteins}

Representing proteins as graphs enables learning over 3D structural relationships:

\subsection{Graph Convolutional Networks (GCN)}

GCNs \cite{kipf2017semi} aggregate neighbor features through spectral convolutions. For proteins, nodes represent residues and edges encode spatial proximity or bonded connectivity.

\subsection{Graph Attention Networks (GAT/GATv2)}

GAT \cite{velickovic2018graph} and GATv2 \cite{brody2022how} learn attention weights between connected nodes, enabling the model to focus on relevant interactions. In protein applications, attention can learn to weight spatially proximal residues differently from distant ones.

\subsection{$k$-Nearest Neighbor Graphs}

Instead of fully-connected or distance-cutoff graphs, $k$-NN graphs connect each residue to its $k$ nearest spatial neighbors. This:
\begin{itemize}
    \item Reduces computational cost from $O(N^2)$ to $O(Nk)$
    \item Forces attention to local structural context
    \item Prevents over-smoothing from global message passing
    \item Empirically improves protein property prediction (our ablation: +0.043 PCC with $k=30$)
\end{itemize}

\section{The MegaScale Dataset}

Tsuboyama et al.\ (2023) \cite{tsuboyama2023mega} measured stability for $\sim$776,000 mutations across $\sim$580 small protein domains using:
\begin{itemize}
    \item Massively parallel proteolysis assay (DMS: Deep Mutational Scanning)
    \item AlphaFold2-predicted structures for all domains
    \item Both $\Delta G$ (absolute stability) and $\Delta\Delta G$ (mutation effect) measurements
\end{itemize}

For this thesis, we use the PNAS-filtered subset:
\begin{itemize}
    \item 368 proteins total (340 train, 28 test)
    \item Single-point mutations only (multi-mutations and insertions/deletions removed)
    \item $\Delta\Delta G$ range: --1.0 to 5.0 kcal/mol (clamped)
    \item AlphaFold PDB structures available on Zenodo (14 MB download)
\end{itemize}

\chapter{Methods}
\label{ch:methods}

\section{DeepEF Architecture Overview}

The Protein Energy Model (PEM) predicts per-residue energy contributions that sum to a protein's total folding energy. The $\Delta\Delta G$ for a mutation is computed as:
\begin{equation}
    \Delta\Delta G = E_{\text{mutant}} - E_{\text{wild-type}} = (E_{\text{unfolded}}^{\text{mut}} - E_{\text{folded}}^{\text{mut}}) - (E_{\text{unfolded}}^{\text{wt}} - E_{\text{folded}}^{\text{wt}})
\end{equation}

\subsection{Input Representation}

For each protein, the input feature vector per residue is:
\begin{equation}
    \mathbf{F}_i = [\mathbf{D}_i \,||\, \mathbf{F}_b \,||\, \mathbf{e}_i \,||\, \mathbf{h}_i] \in \mathbb{R}^{16 + 32 + E + 20}
\end{equation}
where:
\begin{itemize}
    \item $\mathbf{D}_i \in \mathbb{R}^{16}$: Distance matrix features (16 Gaussian kernels applied to atom pair distances, summed over neighbors, L2-normalized)
    \item $\mathbf{F}_b \in \mathbb{R}^{32}$: Bonded features (16 bonded + 16 non-bonded Gaussian distance kernels from immediate sequential neighbors)
    \item $\mathbf{e}_i \in \mathbb{R}^{E}$: Protein language model embedding ($E = 1024$ for ProtT5, $E = 1280$ for SaProt/ESM-2)
    \item $\mathbf{h}_i \in \mathbb{R}^{20}$: One-hot amino acid encoding
\end{itemize}

The distance matrix uses Gaussian radial basis functions:
\begin{equation}
    D_{ij} = \text{ReLU}\left(\exp(\gamma \cdot d_{ij}^2)\right), \quad \gamma = -0.08
\end{equation}
computed between backbone atom pairs (N, CA, C, CB) using 16 distance bins.

\subsection{Dual-Branch GNN}

The model processes protein graphs through two parallel branches:

\textbf{GCN Branch} (bonded + non-bonded context):
\begin{itemize}
    \item Input: $[\mathbf{D}(16) \,||\, \mathbf{F}_b(32) \,||\, \mathbf{h}(20)] = 68$ dimensions (or $52$ without projection)
    \item 3 GCN layers with dropout ($p = 0.2$)
    \item Edge connectivity: fully connected (all residues connected)
    \item Output: 36-dimensional node features
\end{itemize}

\textbf{GAT Branch} (non-bonded attention):
\begin{itemize}
    \item Input: $[\mathbf{D}(16) \,||\, \mathbf{h}(20)] = 36$ dimensions
    \item 3 GATv2 layers with 4 attention heads, dropout ($p = 0.2$)
    \item Edge connectivity: $k$-NN graph ($k = 30$) based on CA atom distances
    \item Output: 36-dimensional node features
\end{itemize}

\subsection{$k$-NN Graph Construction}

For the GAT branch, we construct a directed $k$-nearest neighbor graph:
\begin{equation}
    \mathcal{E} = \{(i, j) : j \in \text{kNN}(i, k=30) \text{ based on } \|\mathbf{r}_i^{CA} - \mathbf{r}_j^{CA}\|_2\}
\end{equation}

This replaces the original fully-connected graph and provides:
\begin{itemize}
    \item Reduced computation: $O(Nk)$ vs $O(N^2)$ edges
    \item Focus on local structural contacts (which determine stability)
    \item Empirically validated: +0.043 PCC improvement over fully-connected
\end{itemize}

\subsection{Feature Fusion and Prediction}

After both GNN branches:
\begin{enumerate}
    \item Concatenate GCN and GAT outputs: $\mathbf{x} = [\mathbf{x}_{\text{GCN}} \,||\, \mathbf{x}_{\text{GAT}}] \in \mathbb{R}^{72}$
    \item Apply instance normalization
    \item Concatenate raw LLM embeddings: $\mathbf{x} = [\mathbf{x} \,||\, \mathbf{e}] \in \mathbb{R}^{72 + E}$
    \item Apply Light Attention mechanism
    \item Two-layer MLP: $(72 + E) \rightarrow 128 \rightarrow 1$
\end{enumerate}

\subsection{Light Attention}

Light Attention \cite{starklight} provides a learnable pooling mechanism over the sequence dimension:
\begin{equation}
    \mathbf{z} = \text{softmax}(\mathbf{W}_a \mathbf{X}^T) \cdot \mathbf{X}
\end{equation}
where $\mathbf{X} \in \mathbb{R}^{N \times d}$ is the sequence of node features and $\mathbf{W}_a$ learns which residues to attend to for the global prediction.

\section{Loss Function}

We use a composite loss combining Huber loss and ranking loss:
\begin{equation}
    \mathcal{L} = \mathcal{L}_{\text{Huber}} + \lambda \cdot \mathcal{L}_{\text{ranking}}
\end{equation}

\textbf{Huber Loss} ($\delta = 1.0$):
\begin{equation}
    \mathcal{L}_{\text{Huber}}(y, \hat{y}) = \begin{cases}
        \frac{1}{2}(y - \hat{y})^2 & \text{if } |y - \hat{y}| \leq \delta \\
        \delta(|y - \hat{y}| - \frac{1}{2}\delta) & \text{otherwise}
    \end{cases}
\end{equation}

\textbf{Ranking Loss} (margin = 0.1, $\lambda = 0.1$):
\begin{equation}
    \mathcal{L}_{\text{ranking}} = \frac{1}{|P|} \sum_{(i,j) \in P} \max(0, -\text{sign}(y_i - y_j)(\hat{y}_i - \hat{y}_j) + m)
\end{equation}
where $P$ is the set of mutation pairs within a protein, encouraging correct relative ordering.

\section{Training Pipeline}

\subsection{Training Configuration}

\begin{table}[h]
\centering
\caption{Training hyperparameters}
\label{tab:hyperparams}
\begin{tabular}{lc}
\toprule
\textbf{Parameter} & \textbf{Value} \\
\midrule
GNN layers & 3 \\
Dropout rate & 0.2 \\
Gaussian coefficient $\gamma$ & --0.08 \\
Batch size & 1 (protein-level) \\
Mini-batch size & 64 (mutations per step) \\
Learning rate & $10^{-4}$ \\
LR schedule & Cosine annealing ($\eta_{\min} = 10^{-6}$) \\
Optimizer & Adam (weight decay $= 10^{-5}$) \\
Gradient accumulation & 4 steps \\
Epochs & 30 \\
$\Delta\Delta G$ clipping & $[-1.0, 5.0]$ kcal/mol \\
\bottomrule
\end{tabular}
\end{table}

\subsection{$\Delta\Delta G$ Computation}

For each protein in a batch:
\begin{enumerate}
    \item Construct folded graph from backbone coordinates + embedding of mutation $m$
    \item Construct unfolded graph (diagonal-only distance matrix, same features)
    \item Predict energies: $E_{\text{folded}}(m)$, $E_{\text{unfolded}}(m)$
    \item $\Delta G(m) = E_{\text{unfolded}}(m) - E_{\text{folded}}(m)$
    \item $\Delta\Delta G(m) = \Delta G(m) - \Delta G(\text{wild-type})$ (index 0 is always WT)
\end{enumerate}

\subsection{From-Scratch Training}

Unlike the original DeepEF which used a pretrained checkpoint (43\_final\_model.pt) from native/decoy discrimination on CASP12, this work trains \textbf{entirely from scratch} on MegaScale mutation data. This means:
\begin{itemize}
    \item No transfer learning from structure discrimination
    \item Model must learn both structural representations AND stability prediction simultaneously
    \item Expected performance ceiling lower than with pretraining ($\sim$0.08--0.12 PCC gap)
\end{itemize}

\section{Data Processing}

\subsection{MegaScale/PNAS Dataset}

The dataset is derived from Tsuboyama et al.\ (2023) with the following filtering:
\begin{itemize}
    \item Remove insertions and deletions (point mutations only)
    \item Remove multi-site mutations (single-point only)
    \item Remove test set homologs from training
    \item Clamp $\Delta\Delta G$ to $[-1.0, 5.0]$ kcal/mol
    \item PNAS mutation list filtering for standardized comparison
\end{itemize}

Final dataset statistics:
\begin{itemize}
    \item \textbf{Training:} 340 proteins, $\sim$50,000--100,000 mutations
    \item \textbf{Test:} 28 proteins (from ThermoMPNN benchmark set)
    \item \textbf{Protein sizes:} 50--200 residues (small domains)
    \item \textbf{Structures:} AlphaFold2-predicted (available on Zenodo)
\end{itemize}

\subsection{Embedding Generation}

Currently: ProtT5 per-residue embeddings (1024-dim) for each mutant sequence, stored as:
\begin{verbatim}
training_data/PROTEIN_NAME/prott5_embeddings/
    prott5_embedding_0.pt   # Chunk 0: [128, seq_len, 1024]
    prott5_embedding_1.pt   # Chunk 1: [128, seq_len, 1024]
    ...
\end{verbatim}

Planned upgrade: SaProt structure-aware embeddings (1280-dim), requiring:
\begin{enumerate}
    \item Download AlphaFold PDB files from Zenodo (14 MB)
    \item Extract 3Di structural tokens via Foldseek
    \item Generate SaProt embeddings with combined AA+3Di input
\end{enumerate}

\chapter{Experiments}
\label{ch:experiments}

This chapter details our systematic experimental process: reproducing the baseline, loss function selection, architectural ablation, failure analysis, and final configuration.

\section{Experimental Setup}

All experiments conducted on:
\begin{itemize}
    \item \textbf{Hardware:} NVIDIA RTX 4070 (8GB VRAM), 32 CPU cores
    \item \textbf{Software:} PyTorch 2.x, PyTorch Geometric, WSL2 (Ubuntu)
    \item \textbf{Training:} From scratch (no pretrained checkpoint)
    \item \textbf{Metric:} Pearson Correlation Coefficient (PCC) on test set (28 proteins)
    \item \textbf{Training time:} $\sim$2--3 hours per experiment (15 epochs)
\end{itemize}

\section{Phase 1: Loss Function Selection}

\subsection{Motivation}
The original DeepEF used L1 loss (Mean Absolute Error). We hypothesized that:
\begin{enumerate}
    \item Huber loss would be more robust to outlier mutations
    \item Ranking loss would improve correlation by preserving relative ordering
\end{enumerate}

\subsection{Experiments}

\begin{table}[h]
\centering
\caption{Loss function ablation results (15 epochs, from scratch)}
\label{tab:loss_ablation}
\begin{tabular}{lc}
\toprule
\textbf{Loss Function} & \textbf{Test PCC} \\
\midrule
L1 (baseline) & 0.467 \\
Huber ($\delta = 1.0$) & 0.478 \\
Ranking only & 0.421 \\
\textbf{Huber + Ranking ($\lambda = 0.1$)} & \textbf{0.486} \\
\bottomrule
\end{tabular}
\end{table}

\subsection{Analysis}

Huber + Ranking achieves the best PCC (0.486) because:
\begin{itemize}
    \item Huber loss reduces sensitivity to outlier $\Delta\Delta G$ values (measurements with high experimental noise)
    \item Ranking loss explicitly optimizes for correlation --- it doesn't care about absolute scale, only relative ordering
    \item The combination balances absolute accuracy (Huber) with correlation (Ranking)
    \item $\lambda = 0.1$ gives ranking a small but consistent push without destabilizing Huber convergence
\end{itemize}

\section{Phase 2: Architecture Ablation}

\subsection{Motivation}
With loss function fixed (Huber + Ranking), we tested three architectural modifications:
\begin{enumerate}
    \item \textbf{Multi-RBF:} Replace single Gaussian kernel with 16 centers (2--22\AA, width=1.5)
    \item \textbf{$k$-NN GAT:} Replace fully-connected GAT edges with $k=30$ nearest neighbors
    \item \textbf{Edge features:} Add 32-dim edge attributes (16 RBF + 16 positional) to GATv2Conv
\end{enumerate}

\subsection{Experimental Design}

Critical lesson learned: \textbf{test one variable at a time}. Our initial attempt (Phase~2 combined) failed catastrophically because multiple changes interacted unpredictably. The corrected ablation:

\begin{table}[h]
\centering
\caption{Architecture ablation results (one variable at a time, 15 epochs)}
\label{tab:arch_ablation}
\begin{tabular}{lccc}
\toprule
\textbf{Experiment} & \textbf{Configuration} & \textbf{Test PCC} & \textbf{$\Delta$ vs Baseline} \\
\midrule
Exp 0 (baseline) & Huber+Ranking, fully-connected & 0.4677 & --- \\
Exp A (Multi-RBF) & + 16-center Gaussian & 0.3871 & \textcolor{red}{--0.081} \\
\textbf{Exp B ($k$-NN GAT)} & \textbf{+ $k$=30 nearest neighbors} & \textbf{0.5103} & \textcolor{green!50!black}{\textbf{+0.043}} \\
Exp C (Edge features) & + 32-dim edge attributes & 0.4643 & --0.003 \\
Exp D ($k$-NN + Edge) & B + C combined & 0.4479 & --0.020 \\
Exp E (All three) & A + B + C combined & 0.4208 & --0.047 \\
\bottomrule
\end{tabular}
\end{table}

\subsection{Key Findings}

\begin{enumerate}
    \item \textbf{$k$-NN GAT is the only beneficial change (+0.043 PCC).} Constraining attention to local neighbors forces the model to learn from physically relevant contacts.

    \item \textbf{Multi-RBF hurts significantly (--0.081 PCC).} The 16-center expansion likely overfits to distance distributions in training proteins without generalizing.

    \item \textbf{Edge features are neutral (--0.003 PCC).} GATv2's learned attention already captures distance information implicitly.

    \item \textbf{Combinations are worse than individual changes.} Multi-RBF's negative effect dominates all combinations. Even $k$-NN + Edge (without Multi-RBF) underperforms $k$-NN alone.
\end{enumerate}

\section{Phase 2 Failure: Lessons Learned}

\subsection{The Failure}

Before the controlled ablation, we attempted to deploy all three changes simultaneously (``Phase 2 bundle''). Results were catastrophic:
\begin{itemize}
    \item Path A (Phase 1 extended): PCC = 0.450 (regression from 0.486)
    \item Path B (Phase 2 architecture): PCC = 0.086--0.179 (near-random)
\end{itemize}

\subsection{Root Cause Analysis}

Five bugs were identified through code analysis:

\begin{enumerate}
    \item \textbf{Bug 1: Default parameter poisoning.} The \texttt{use\_multi\_rbf=True} default in \texttt{train\_utils.py:get\_graph()} affected ALL experiments, including the ``baseline'' Path A. Every experiment unknowingly used Multi-RBF.

    \item \textbf{Bug 2: Dead code path.} \texttt{pnas\_train.py} calls the model WITHOUT \texttt{ca\_coords}, so the $k$-NN graph and edge features were never actually computed. The model fell back to fully-connected with edge\_dim=32 but no edge data.

    \item \textbf{Bug 3: GATv2Conv initialization mismatch.} GATv2Conv was initialized with \texttt{edge\_dim=32} (expecting edge features) but received \texttt{edge\_attr=None} during forward pass, causing silent interference.

    \item \textbf{Bug 4: Embedding projection never activated.} The \texttt{emb\_projection} parameter defaults to \texttt{"none"} and was never passed from \texttt{pnas\_train.py} --- ruled out as cause but initially suspected.

    \item \textbf{Bug 5: Loss function not in main branch.} Huber + Ranking loss was only on the collaborator's local branch, not the git main branch code.
\end{enumerate}

\subsection{Lessons}

\begin{enumerate}
    \item \textbf{Never bundle multiple changes.} Each experiment must change exactly one variable.
    \item \textbf{Use CLI flags, not config file defaults.} Explicit $>$ implicit.
    \item \textbf{Verify code paths actually execute.} Add logging/assertions to confirm features are active.
    \item \textbf{Read the data flow end-to-end} before concluding what code actually runs.
\end{enumerate}

\section{Final Best Configuration}

After the corrected ablation, our best configuration:

\begin{table}[h]
\centering
\caption{Current best model configuration}
\label{tab:best_config}
\begin{tabular}{ll}
\toprule
\textbf{Component} & \textbf{Setting} \\
\midrule
Loss function & Huber ($\delta=1.0$) + Ranking ($\lambda=0.1$, margin=0.1) \\
GAT edge construction & $k$-NN ($k=30$) based on CA distances \\
GCN edge construction & Fully connected \\
GNN layers & 3 (both branches) \\
Embeddings & ProtT5 (1024-dim) \\
Training & From scratch, 15 epochs, LR=$10^{-4}$ \\
\textbf{Result} & \textbf{PCC = 0.510} \\
\bottomrule
\end{tabular}
\end{table}

\chapter{Current Results}
\label{ch:results}

\section{Summary of Progress}

\begin{table}[h]
\centering
\caption{Chronological progress from baseline to current best}
\label{tab:progress}
\begin{tabular}{clcc}
\toprule
\textbf{Step} & \textbf{Change} & \textbf{PCC} & \textbf{Gain} \\
\midrule
0 & L1 loss, fully-connected, from scratch & 0.467 & --- \\
1 & + Huber + Ranking loss & 0.486 & +0.019 \\
2 & + $k$-NN GAT ($k=30$) & \textbf{0.510} & +0.024 \\
\midrule
\multicolumn{2}{l}{\textbf{Total improvement}} & & \textbf{+0.043} \\
\bottomrule
\end{tabular}
\end{table}

\section{Comparison with Related Work}

\begin{table}[h]
\centering
\caption{Comparison with published methods on MegaScale/PNAS benchmark}
\label{tab:comparison}
\begin{tabular}{lcc}
\toprule
\textbf{Method} & \textbf{PCC (approx.)} & \textbf{Uses Structure?} \\
\midrule
BLOSUM62 baseline & $\sim$0.30 & No \\
ESM-2 zero-shot (log-likelihood ratio) & $\sim$0.45 & No \\
DeepEF (this work, current) & 0.510 & Yes (coordinates) \\
DeepEF (original, with pretraining) & $\sim$0.54 & Yes \\
RaSP & $\sim$0.60 & Yes \\
ThermoMPNN & $\sim$0.75 & Yes \\
\bottomrule
\end{tabular}
\end{table}

Note: Comparisons are approximate as methods use different test splits and evaluation protocols.

\section{What We Have vs.\ What We Need}

\begin{table}[h]
\centering
\caption{Current state assessment}
\label{tab:state}
\begin{tabular}{p{6cm}p{6cm}}
\toprule
\textbf{What We Have (Completed)} & \textbf{What We Need (Remaining)} \\
\midrule
\vspace{0.1cm}
$\checkmark$ Working training pipeline (from scratch) \newline
$\checkmark$ Huber + Ranking loss (validated) \newline
$\checkmark$ $k$-NN GAT architecture ($k=30$) \newline
$\checkmark$ Controlled ablation methodology \newline
$\checkmark$ PCC = 0.510 on 28-protein test set \newline
$\checkmark$ MegaScale data preprocessed \newline
$\checkmark$ AlphaFold PDB files identified (Zenodo) \newline
$\checkmark$ Code supports CLI flags for all experiments
&
\vspace{0.1cm}
$\square$ Training improvements (cosine LR, grad accum) \newline
$\square$ SaProt embeddings (1280-dim, structure-aware) \newline
$\square$ Download PDBs + install Foldseek \newline
$\square$ Per-protein fine-tuning on test proteins \newline
$\square$ Multi-seed ensemble (5 seeds) \newline
$\square$ Embedding projection into GNN branches \newline
$\square$ Final PCC target: $\geq$ 0.70
\\
\bottomrule
\end{tabular}
\end{table}

\section{Analysis of Current Limitations}

\subsection{Gap to Target}

Current PCC = 0.510, target $\geq$ 0.70. The gap of +0.19 must come from:
\begin{enumerate}
    \item \textbf{Input quality:} ProtT5 embeddings are sequence-only (Spearman $\sim$0.65). SaProt is structure-aware (Spearman 0.724). This single change should provide +0.03--0.05 PCC.
    \item \textbf{Training suboptimality:} Only 15 epochs with constant LR. Cosine schedule + more epochs + gradient accumulation = smoother convergence.
    \item \textbf{No pretraining:} The original DeepEF used native/decoy pretraining. Without it, the model must learn both structural representations and stability prediction simultaneously.
    \item \textbf{Single model variance:} A single random seed has high variance. Ensembling 5 seeds reduces noise by $\sim\sqrt{5}$.
\end{enumerate}

\subsection{What Cannot Be Fixed}

\begin{itemize}
    \item \textbf{No pretrained checkpoint:} The 43\_final\_model.pt from CASP12 pretraining is unavailable. This likely accounts for $\sim$0.08--0.12 PCC gap.
    \item \textbf{Dataset size:} 340 training proteins with $\sim$50K mutations is sufficient but not as large as ThermoMPNN's training set.
    \item \textbf{Hardware constraint:} 8GB VRAM limits batch size and model scale.
\end{itemize}

\chapter{Planned Improvements and Future Work}
\label{ch:future_work}

\section{Improvement Roadmap}

We present a four-stage plan to systematically close the gap from PCC = 0.510 to $\geq$ 0.70. Each stage addresses a different bottleneck: underfitting $\rightarrow$ input quality $\rightarrow$ task specificity $\rightarrow$ variance reduction.

\begin{table}[h]
\centering
\caption{Four-stage improvement plan with expected gains}
\label{tab:roadmap}
\begin{tabular}{clccc}
\toprule
\textbf{Stage} & \textbf{Improvement} & \textbf{Expected PCC} & \textbf{P(helps)} & \textbf{Time} \\
\midrule
1 & Training optimization & 0.53--0.55 & 85\% & 1 day \\
2 & SaProt embeddings & 0.56--0.60 & 70\% & 1--2 days \\
3 & Per-protein fine-tuning & 0.59--0.65 & 55\% & 1 day \\
4 & 5-seed ensemble & 0.62--0.70 & 90\% & 1 day \\
\bottomrule
\end{tabular}
\end{table}

\section{Stage 1: Training Optimization}

\textbf{Changes:}
\begin{itemize}
    \item Cosine annealing LR schedule: $10^{-4} \rightarrow 10^{-6}$ over 30 epochs
    \item Increase training from 15 to 30 epochs
    \item Gradient accumulation: effective batch size = 4
    \item Weight decay: $10^{-5}$ (L2 regularization)
\end{itemize}

\textbf{Rationale:}
\begin{itemize}
    \item Constant LR causes oscillation in later epochs instead of fine convergence
    \item 15 epochs may not be sufficient for the model to converge
    \item Batch size = 1 produces noisy gradients; accumulation smooths optimization
    \item 22M parameters with effective batch=1 risks overfitting; weight decay regularizes
\end{itemize}

\textbf{Risk:} None. These are standard practices with no downside.

\section{Stage 2: SaProt Embeddings}

\textbf{Change:} Replace ProtT5 (1024-dim, sequence-only) with SaProt (1280-dim, structure-aware).

\textbf{Why SaProt:}
\begin{itemize}
    \item SaProt encodes both amino acid identity AND local backbone geometry (via Foldseek 3Di tokens)
    \item Published thermostability Spearman: 0.724 (vs ProtT5 $\sim$0.65, ESM-2 = 0.680)
    \item $\Delta\Delta G$ prediction depends on 3D context (burial, contacts) --- exactly what SaProt provides
    \item Same architecture as ESM-2 (650M params) but with structure-aware vocabulary
\end{itemize}

\textbf{Implementation requirements:}
\begin{enumerate}
    \item Download AlphaFold PDB files from Zenodo (14 MB) --- URL already in codebase
    \item Install Foldseek binary (structural alphabet extraction)
    \item Run \texttt{structureto3didescriptor} to get 3Di tokens per protein
    \item Generate per-mutation SaProt embeddings ($\sim$2--4 hours on GPU)
    \item Update \texttt{model\_cfg.py}: \texttt{emb\_input\_dim = 1280}
    \item Fix hardcoded index: \texttt{self.llm\_index = -(emb\_input\_dim + 20)}
\end{enumerate}

\textbf{Risk:} Low. SaProt is a published model (ICLR 2024), MIT license, 54K+ downloads. PDB files confirmed available.

\section{Stage 3: Per-Protein Fine-Tuning}

\textbf{Change:} After training the global model, fine-tune separate copies on each test protein's own mutation data.

\textbf{Rationale:}
\begin{itemize}
    \item The global model learns average $\Delta\Delta G$ patterns across 340 diverse proteins
    \item Each protein has a unique energy landscape (different fold, different stability determinants)
    \item Per-protein adaptation captures protein-specific effects (e.g., a tight hydrophobic core vs.\ a flexible loop protein)
    \item This approach is used by ThermoMPNN and RaSP for their best reported numbers
\end{itemize}

\textbf{Implementation:}
\begin{enumerate}
    \item Load best global model checkpoint
    \item For each test protein: fine-tune for 5--10 epochs on that protein's training mutations
    \item Predict on held-out mutations
    \item Report per-protein and aggregate PCC
\end{enumerate}

\textbf{Risk:} Medium. Helps for proteins with many mutations, may overfit for proteins with few ($<$20) mutations.

\section{Stage 4: Ensemble}

\textbf{Change:} Train 5 models with different random seeds (42, 123, 456, 789, 1337), average predictions.

\textbf{Rationale:}
\begin{itemize}
    \item Each random seed finds a different local minimum
    \item Averaging reduces prediction variance by factor of $\sim\sqrt{5}$
    \item Empirically always improves correlation metrics (+0.02--0.03 PCC typical)
    \item Zero implementation complexity --- just run 5 times
\end{itemize}

\textbf{Risk:} None. Mathematically guaranteed to reduce variance. Only cost is 5$\times$ training time.

\section{Additional Directions Under Consideration}

\subsection{Embedding Projection into GNN (emb\_projection=``mlp'')}

The model already contains code to project the LLM embedding (1280-dim) $\rightarrow$ 16-dim and feed it directly into the GCN and GAT branches. Currently disabled (\texttt{emb\_projection="none"}). This would give the GNN access to semantic information during message passing, not just at the output stage.

\textbf{P(helps): 55\% | Expected gain: +0.01--0.03 PCC}

\subsection{Embedding-Enriched One-Hot}

Project embedding $\rightarrow$ 20-dim and ADD to the one-hot vector, giving the GNN amino acid identity enriched with evolutionary/structural context.

\textbf{P(helps): 50\% | Expected gain: +0.01--0.03 PCC}

\subsection{GCN Branch Modification}

The GCN branch uses fully-connected edges (every residue connected). Applying $k$-NN to GCN as well (or removing GCN entirely since GAT dominates) could reduce noise.

\textbf{P(helps): 40\% | Expected gain: +0.01--0.02 PCC}

\section{Decision Tree}

After each stage, the next action depends on the result:

\begin{verbatim}
After Stage 1 (PCC = ___):
  >= 0.53  --> Continue to Stage 2
  <  0.52  --> Debug (try without cosine LR)

After Stage 2 (PCC = ___):
  >= 0.63  --> Skip Stage 3, go to Stage 4 (ensemble)
  0.55-0.63 --> Continue to Stage 3
  <  0.55  --> Try dual embeddings (ProtT5 + SaProt concat)

After Stage 3 (PCC = ___):
  >= 0.68  --> Ensemble will reach 0.70+
  0.62-0.68 --> Ensemble --> 0.65-0.72

After Stage 4 (PCC = ___):
  >= 0.70  --> TARGET REACHED
  0.65-0.70 --> Strong thesis result
  <  0.65  --> Data/pretraining ceiling
\end{verbatim}

\section{Realistic Expectations}

\begin{table}[h]
\centering
\caption{Probability assessment for final PCC}
\label{tab:probability}
\begin{tabular}{lc}
\toprule
\textbf{Final PCC Range} & \textbf{Probability} \\
\midrule
$\geq$ 0.75 & 20--30\% \\
0.70--0.75 & 25--35\% \\
0.65--0.70 & 30--35\% \\
$<$ 0.65 & 10--20\% \\
\bottomrule
\end{tabular}
\end{table}

\textbf{Key constraint:} Without the pretrained checkpoint (native/decoy discrimination), we estimate a $\sim$0.08--0.12 PCC ceiling reduction. ThermoMPNN achieves $\sim$0.75 with structure-specific pretraining AND larger training sets. Our realistic ceiling is 0.70--0.72 with all improvements applied.

\section{Timeline}

\begin{table}[h]
\centering
\caption{Estimated timeline for remaining work}
\label{tab:timeline}
\begin{tabular}{lll}
\toprule
\textbf{Week} & \textbf{Task} & \textbf{Deliverable} \\
\midrule
Week 1 & Stage 1 (training) + download PDBs & PCC after Stage 1 \\
Week 1--2 & Generate SaProt embeddings & Embeddings ready \\
Week 2 & Stage 2 (SaProt training) & PCC after Stage 2 \\
Week 2--3 & Stage 3 (per-protein FT) & PCC after Stage 3 \\
Week 3 & Stage 4 (ensemble) & Final PCC \\
Week 3--4 & Analysis, writing, figures & Complete thesis \\
\bottomrule
\end{tabular}
\end{table}


\bibliographystyle{plain}
\bibliography{references}

\end{document}
